% TOP_PROBLEM: {"id": 10000049, "title": "Liouville extensions by an isometric integer action", "category_id": 3, "category": {"id": 3, "name": "graph_theory", "display_name": "Graph Theory", "description": "Problems involving graphs, networks, and their properties.", "slug": "graph-theory", "order_index": 3, "created_at": "2026-07-31T15:26:25.671Z"}, "status": "open", "problem_number": "AMR-099-0049", "set_id": 13, "statusQualification": "Makes explicit the quotient walk inherited from simple random walk, preserving edge multiplicities and holding probabilities."}
% FIRST_POSTED: 2026-10-01
% ENTRY_KIND: statement_only
% UnsolvedMath ID 10000049; source catalogue code AMR-099-0049
% !TeX program = lualatex
% Definition and sources only; this file is not an LLM solution attempt.
\documentclass[11pt,a4paper]{article}
\usepackage[margin=25mm]{geometry}
\usepackage{fontspec}
\setmainfont{lmroman10-regular.otf}[BoldFont=lmroman10-bold.otf,
 ItalicFont=lmroman10-italic.otf,BoldItalicFont=lmroman10-bolditalic.otf]
\usepackage{amsmath,amssymb,mathtools}
\usepackage{unicode-math}
\setmathfont{latinmodern-math.otf}
\AtBeginDocument{\let\setminus\smallsetminus}
\usepackage{xurl,hyperref}
\hypersetup{colorlinks=true,urlcolor=blue}
\setcounter{secnumdepth}{0}
\setlength{\parindent}{0pt}
\setlength{\parskip}{5pt}
\setlength{\emergencystretch}{3em}
\begin{document}
\section*{Liouville extensions by an isometric integer action}\label{10000049}
{\small UnsolvedMath ID 10000049 · AMR-099-0049 · Graph Theory · AMR Open Problem Lists}
\subsection{Definitions and mathematical statement}
Let $G$ be a connected locally finite graph, and let an automorphism $T$ define an action of $\mathbb Z$ by $n\cdot v=T^nv$. Let $\pi:V(G)\to V(G)/\mathbb Z$ be the orbit map and choose a set $F$ containing exactly one representative of each orbit. Its translates are $T^kF$, $k\in\mathbb Z$.

The quotient random walk is the projection of simple random walk on $G$. Its transition kernel is
\[
\overline P([v],[w])=
\frac{|\{u\sim_Gv:\pi(u)=[w]\}|}{\deg_G(v)}.
\]
This does not depend on the representative $v$, because $T$ preserves adjacency. Retaining these multiplicities and possible holding probabilities specifies the quotient convention. Suppose the quotient is Liouville: every bounded real $g$ satisfying $\overline Pg=g$ is constant.

Assume that for the chosen $F$, simple random walk $(X_n)$ on $G$, from every starting vertex, satisfies
\[
\mathbb P\bigl(X_n\in T^kF\text{ for infinitely many }n\bigr)=1
\qquad\text{for every }k\in\mathbb Z.
\]
Must every bounded harmonic function on $G$ itself be constant?

The recurrence assumption is recurrence to each translated set, which may be infinite; it does not assert recurrence to individual vertices. Countability allows the visits to all translated fundamental sets to hold simultaneously almost surely. The action need not be transitive on $G$. The conclusion asks whether bounded harmonic information can remain in the direction of the integer action when the quotient carries none and the walk repeatedly returns to every translate of a fundamental domain.
\subsection{Short English statement}
Does a Liouville quotient and repeated returns to all translated fundamental domains force an integer extension of a graph to be Liouville?
\subsection{Sources}
\begin{itemize}
\item[S1] ULAM.AI and the UnsolvedMath Contributors, supplied problems.json, record 10000049 (AMR-099-0049), AMR Open Problem Lists. Catalogue identity and source transcription; CC BY 4.0.
\url{https://huggingface.co/datasets/ulamai/UnsolvedMath}.
\item[S2] Benjamini - Coarse Geometry and Randomness (Saint-Flour notes), Open problem 13.3, PDF page 107. Original source indexed by AMR-099-0049.
\url{https://www.wisdom.weizmann.ac.il/~itai/stflouraug24.pdf#page=107}.
\item[S3] I. Benjamini, Coarse Geometry and Randomness, primary Saint-Flour notes; the numbered question and its surrounding definitions.
\url{https://arquivo.pt/noFrame/replay/20201231041548id_/http://www.wisdom.weizmann.ac.il/~itai/stflouraug24.pdf}.
\end{itemize}
\end{document}
